% TOP_PROBLEM: {"id": 3352, "title": "Twin prime conjecture", "category_id": 1, "category": {"id": 1, "name": "number_theory", "display_name": "Number Theory", "description": "Properties of integers, prime numbers, Diophantine equations.", "slug": "number-theory", "order_index": 1, "created_at": "2026-07-31T15:26:25.670Z"}, "status": "open", "problem_number": "OPG-36952", "set_id": 12}
% FIRST_POSTED: 2026-10-01
% ENTRY_KIND: statement_only
% UnsolvedMath ID 3352; source catalogue code OPG-36952
% !TeX program = lualatex
% Definition and sources only; this file is not an LLM solution attempt.
\documentclass[11pt,a4paper]{article}
\usepackage[margin=25mm]{geometry}
\usepackage{fontspec}
\setmainfont{lmroman10-regular.otf}[BoldFont=lmroman10-bold.otf,
 ItalicFont=lmroman10-italic.otf,BoldItalicFont=lmroman10-bolditalic.otf]
\usepackage{amsmath,amssymb,mathtools}
\usepackage{unicode-math}
\setmathfont{latinmodern-math.otf}
\AtBeginDocument{\let\setminus\smallsetminus}
\usepackage{xurl,hyperref}
\hypersetup{colorlinks=true,urlcolor=blue}
\setcounter{secnumdepth}{0}
\setlength{\parindent}{0pt}
\setlength{\parskip}{5pt}
\setlength{\emergencystretch}{3em}
\begin{document}
\section*{Twin prime conjecture}\label{3352}
{\small UnsolvedMath ID 3352 · OPG-36952 · Number Theory · OpenGarden}
\subsection{Definitions and mathematical statement}
A prime is an integer $p>1$ whose only positive divisors are $1$ and $p$.
A twin-prime pair is an unordered pair $\{p,p+2\}$ for which both
entries are prime. Define its counting function by
\[
 \pi_2(x)=\#\{p\in\mathbb Z:2\le p\le x,
                           \ p\text{ and }p+2\text{ are prime}\}.
\]
Here the cutoff is on the smaller member; choosing the larger member
instead would not change the infinitude question. The twin-prime
conjecture asserts that $\pi_2(x)\to\infty$ as $x\to\infty$.
Equivalently, for every real $X$ there exists an integer $p>X$ such
that $p$ and $p+2$ are prime.

The common difference is fixed at exactly two before $X$ varies. Thus
allowing a bounded but unspecified even gap changes the target.
No asymptotic formula, positive density among primes, or effective
upper bound for the next pair is part of this statement. Except for
the pair $\{3,5\}$, a twin pair has the form $\{6m-1,6m+1\}$; this
necessary residue condition does not assert primality. A resolution
must establish infinitude for the exact gap or show that only finitely
many such pairs exist.
\subsection{Short English statement}
Are there prime pairs separated by exactly two beyond every numerical bound?
\subsection{Sources}
\begin{itemize}
\item[S1] ULAM.AI and the UnsolvedMath Contributors, supplied problems.json, record 3352 (OPG-36952), OpenGarden collection. Catalogue identity and source transcription; CC BY 4.0.
\url{https://huggingface.co/datasets/ulamai/UnsolvedMath}.
\item[S2] Open Problem Garden, Twin prime conjecture. Original problem and discussion; the mathematical formulation below is expanded from the supplied source record.
\url{http://www.openproblemgarden.org/op/twin_prime_conjecture}.
\end{itemize}
\end{document}
